\documentclass{article}
\usepackage[final]{neurips_2024}
\usepackage[utf8]{inputenc}
\usepackage[T1]{fontenc}
\usepackage{hyperref}
\usepackage{url}
\usepackage{booktabs}
\usepackage{amsfonts}
\usepackage{nicefrac}
\usepackage{microtype}
\usepackage{graphicx}
\usepackage{amsmath}
\usepackage{amssymb}
\usepackage{multirow}
\usepackage{xcolor}
\usepackage{pifont}
\usepackage{algorithm}
\usepackage{algpseudocode}

\title{Beyond Token Space: An Information Augmentation Spectrum\\for Masked Diffusion Language Models}

\author{
  Anonymous Authors
}

\begin{document}
\maketitle

\begin{abstract}
Masked diffusion language models (MDLMs) generate text via iterative denoising, yet each step discards all continuous representations---logits, attention patterns, hidden states---before the next step begins.
This ``Information Island'' problem limits inference-time reasoning.
We introduce the \textbf{Information Augmentation Spectrum}, a four-level hierarchy that systematically adds cross-step information at increasing computational cost:
(1)~Vanilla (no transfer),
(2)~DMI---Diffusion Memory Injection (embedding-level),
(3)~SCP---Selective Confidence Propagation (token-level), and
(4)~DTA---Denoising-Time Adaptation (parameter-level via online LoRA updates).
Full-scale evaluation on Countdown-500 with Dream-7B reveals a surprising finding: \textbf{shallower interventions outperform deeper ones}.
DMI achieves 9.3\% accuracy (${\sim}2{\times}$ vanilla, $p < 0.05$) at near-zero cost ($1.2{\times}$ overhead);
SCP reaches 9.1\% but at $12{\times}$ overhead;
DTA achieves only 4.8\% ($\approx$ vanilla).
Existing remasking methods are equally ineffective: ReMDM-conf scores 4.4\%, with token-level diagnostics revealing only 31.3\% correction precision.
These results demonstrate that effective cross-step information transfer for DLMs operates at the embedding level, not in token or parameter space, challenging the assumption that richer representations yield better inference-time scaling.
\end{abstract}

%% ============================================================
\section{Introduction}
\label{sec:intro}

Masked diffusion language models (MDLMs) generate text through iterative denoising: starting from a fully masked sequence, the model progressively reveals tokens over $T$ steps, producing a complete generation at step $t = 0$ \citep{sahoo2024mdlm, nie2025llada, dream2025}. This iterative structure grants MDLMs unique advantages over autoregressive (AR) models---parallel token prediction, arbitrary generation order, and natural support for self-correction via remasking \citep{wang2025remdm, core2026}. Recent scaling efforts have demonstrated concrete gains: Dream~7B achieves 16.0 on the Countdown planning benchmark versus 6.2 for its AR counterpart \citep{dream2025}, and methods such as remasking-based refinement (ReMDM; \citealp{wang2025remdm}), context probing (CORE; \citealp{core2026}), and reward-guided correction (MDPO+RCR; \citealp{mdpo2025}) leverage the denoising loop to improve generation quality at inference time.

Yet a fundamental limitation persists across all existing DLM inference-time scaling methods: \textbf{each denoising step operates as an information island}. At step $t$, the model receives the current partially masked sequence $\mathbf{y}^{(t)}$, computes a full forward pass to predict masked positions, samples tokens, and then \emph{discards} all intermediate representations---logits, attention patterns, hidden states---before proceeding to step $t{-}1$. The next step starts from scratch with only the discrete token sequence, having no access to the continuous reasoning state accumulated during the previous step. We term this the \emph{cross-step information loss problem}; \citet{xia2026metastate} independently identify it as the ``Information Island'' problem and propose MetaState, a trained GRU+CrossAttention module to bridge denoising steps---but this requires additional training and architectural modification.

Current inference-time scaling approaches attempt to compensate for this cross-step information loss, but they operate exclusively in \emph{token space}: remasking methods (ReMDM-conf, RCR) selectively re-mask and re-predict low-confidence tokens; context-probing methods (CORE) perturb the masked context to detect fragile tokens. These token-level interventions face an inherent signal quality problem. Our full-scale experiments on Countdown-500 with Dream-7B reveal that ReMDM-conf achieves only 4.4\% accuracy (versus 4.7\% for vanilla denoising), and RCR reaches just 5.7\%---marginal or no improvement. Token-level diagnostics expose the root cause: ReMDM-conf has a correction precision of only 31.3\%, meaning it remasks ${\sim}70$\% correct tokens, and creates 94.8 unstable token positions per sample (${\sim}37$\% of the generation area). Remasking without cross-step memory is akin to an editor who can re-read an essay but has lost all notes, margin comments, and reasoning from the previous editing pass---it produces token churning, not genuine refinement.

\textbf{Our key insight} is that MDLM denoising is \emph{structurally analogous to test-time training} (TTT) \citep{sun2024ttt, titans2025}. At each denoising step, the model performs masked language modeling on a progressively revealed sequence---precisely the self-supervised objective that TTT methods use to adapt model parameters at inference time. But whereas AR-based TTT iterates along \emph{sequence positions} (requiring causal structure), DLM denoising iterates along \emph{time steps}, with each step providing full bidirectional context over the entire sequence. This structural alignment means that the denoising objective is \emph{inherently} self-supervised: no auxiliary loss or external supervision is needed. Current DLMs simply never exploit this opportunity---they perform implicit TTT (masked prediction at each step) without ever updating parameters.

We propose \textbf{Denoising-Time Adaptation (DTA)}, a training-free method that converts this implicit TTT into explicit test-time learning. DTA augments standard DLM denoising with a lightweight parameter update at each step: after the model reveals new tokens (E-step), it masks a subset of the revealed tokens, computes the masked language modeling loss, and performs a single gradient step on a zero-initialized LoRA adapter (M-step). The adapter accumulates information across denoising steps, building a \emph{parameter-level memory} of the emerging generation's patterns and constraints (see Figure~\ref{fig:dta_overview} for a visual overview). We formalize this intuition through the \textbf{Variational Denoising-Time Adaptation (VDTA)} framework, which interprets DTA as expectation-maximization in the joint space of tokens $\mathbf{x}_0$ and adapter parameters $\Delta\theta$. We prove that each E-M step improves the variational lower bound (under mild regularity conditions detailed in Section~\ref{sec:vdta}) and provide empirical evidence that the adapter's information about the target sequence increases across denoising steps (Section~\ref{sec:degradation}).

To systematically ablate the value of cross-step information transfer, we introduce the \textbf{Information Augmentation Spectrum}, a hierarchy of four methods that add progressively richer forms of cross-step information---from embedding injection (DMI) to token-level error detection (SCP) to parameter-level adaptation (DTA)---at increasing computational cost and expressivity. This spectrum provides a principled ablation framework: each level adds a qualitatively different form of cross-step memory, enabling us to isolate the contribution of embedding-level, token-level, and parameter-level information transfer (detailed in Section~\ref{sec:spectrum}).

We evaluate DTA and the full information augmentation spectrum on three reasoning benchmarks---Countdown-500, GSM8K-1319, and MBPP-500---using Dream-7B-Instruct as the primary model and LLaDA-8B-Instruct for cross-model verification. All experiments use 3 random seeds with McNemar tests and Bonferroni-corrected bootstrap 95\% confidence intervals. Our contributions are:

\begin{enumerate}
    \item \textbf{Denoising-Time Adaptation (DTA)}: A training-free method that adds online LoRA updates during DLM denoising, establishing the first connection between test-time training and masked diffusion language models. DTA is the only method satisfying ``training-free $+$ parameter-level memory $+$ theoretical guarantee'' simultaneously. On Countdown-500, our simplest spectrum method DMI achieves 9.3\% (${\sim}2{\times}$ vanilla, $p < 0.05$) at near-zero cost; pilot results show DTA's strongest gains on code generation (MBPP: +12.5pp).
    \item \textbf{VDTA theoretical framework}: A variational interpretation proving ELBO monotonicity and formalizing DLM denoising as implicit test-time training, with empirical validation of information accumulation across denoising steps.
    \item \textbf{Information Augmentation Spectrum}: A principled ablation framework (Vanilla $<$ DMI $<$ SCP $<$ DTA) that systematically quantifies the value of cross-step information transfer at increasing computational cost.
    \item \textbf{Comprehensive empirical analysis}: Full-scale evaluation on three benchmarks with two DLM architectures, token-level diagnostic metrics (Correction Precision/Recall, trajectory stability), and ablation studies revealing the memory--stability tradeoff controlled by the decay factor $\gamma$.
\end{enumerate}

%% ============================================================
\section{Background and Related Work}
\label{sec:related}

Denoising-Time Adaptation sits at the intersection of masked diffusion language models, inference-time scaling, and test-time training. We review each thread and highlight the specific gap that DTA fills.

\subsection{Masked Diffusion Language Models}

Discrete diffusion for text generation has progressed rapidly from proof-of-concept to industrial scale. MDLM \citep{sahoo2024mdlm} introduced a Rao-Blackwellized objective that closed the perplexity gap with autoregressive models; SEDD \citep{lou2024sedd} extended score matching to discrete spaces. These foundations enabled two landmark systems: \textbf{LLaDA~8B} \citep{nie2025llada}, which matches LLaMA3~8B on in-context learning and breaks the reversal curse, and \textbf{Dream~7B} \citep{dream2025}, which initializes from an AR checkpoint and introduces context-adaptive noise rescheduling, achieving state-of-the-art planning performance (Countdown: 16.0 vs.\ AR 6.2). The ecosystem has since expanded rapidly \citep{llada15_2025, llada_moe_2025, dreamcoder2025, geminidiffusion2025, svete2025reasoning, jiang2025optimal}.

A defining property of the DLM denoising process is that each step $t$ receives a partially masked sequence $\mathbf{x}_t$, makes predictions via a full bidirectional forward pass, samples tokens, and then \textbf{discards} all continuous representations---logits, attention maps, hidden states---before proceeding to step $t{-}1$. We term this the cross-step information loss problem; \citet{xia2026metastate} independently identify it as the ``Information Island'' problem. Every denoising step starts from scratch, unable to leverage the understanding accumulated by its predecessors.

\subsection{Inference-Time Scaling for DLMs}

We organize existing methods by the mechanism through which they inject cross-step information.

\paragraph{Token-space methods.} ReMDM \citep{arriola2025remdm} introduced four remasking strategies derived from custom backward processes. Running Confidence Remasking (RCR) \citep{chen2025mdpo} tracks per-token confidence across steps. CORE \citep{zhai2026core} probes context robustness by masking neighboring tokens, yielding +9.2pp on MBPP with LLaDA-8B. HEX \citep{lee2025hex} ensembles multiple denoising schedules. Search-based methods---Prism \citep{wang2026prism}, UnMaskFork \citep{li2026unmaskfork}, TReASURe \citep{chen2025treasure}, Reward-Guided Stitching \citep{stitching2026}, PG-DLM \citep{pgdlm2025}, Self-Rewarding SMC \citep{srsmc2026}---achieve strong results but depend on external verifiers, reward models, or multiple parallel trajectories. A common limitation is that they operate exclusively in discrete token space: they reshuffle which tokens are masked and re-predicted but \textbf{never update the model's internal representation}.

\paragraph{Training-dependent methods.} RemeDi \citep{xu2025remedi} adds a dual-stream architecture trained via SFT and RL. ProSeCo \citep{proseco2026} trains a corrector network. RL-based approaches---d1/diffu-GRPO \citep{d12025}, wd1 \citep{wd12025}, DCoLT \citep{dcolt2025}, MDPO \citep{chen2025mdpo}, AGRPO \citep{agrpo2025}---frame denoising as a sequential decision process. MetaState \citep{xia2026metastate} trains a GRU + cross-attention module to maintain persistent hidden state across denoising steps.

\paragraph{Gradient-guided inference.} TABES \citep{tabes2026} derives a Token Influence Score from a first-order expansion of a trajectory cost functional. ETS \citep{ets2026} factorizes the transition probability into a reference policy plus an energy term. Both use gradients at inference time, but they guide \textbf{token selection} rather than updating model parameters. By ``parameter update at test time,'' we mean persistent modification of model weights that carry forward to subsequent denoising steps, as opposed to one-shot gradient signals used to score tokens at a single step.

\paragraph{Positioning summary.} Table~\ref{tab:positioning} situates DTA within this landscape. DTA is the only method that simultaneously satisfies five desiderata: (i) parameter-level updates at test time, (ii) persistent cross-step memory, (iii) no training requirement, (iv) no external verifier, and (v) a theoretical guarantee (ELBO monotonicity; see Section~\ref{sec:vdta}).

\begin{table}[t]
\centering
\caption{Method positioning along key capability axes. DTA uniquely satisfies all five criteria.}
\label{tab:positioning}
\small
\begin{tabular}{lccccc}
\toprule
\textbf{Method} & \textbf{Param.\ update} & \textbf{Cross-step} & \textbf{Training-} & \textbf{No external} & \textbf{Theoretical} \\
 & \textbf{at test time} & \textbf{memory} & \textbf{free} & \textbf{verifier} & \textbf{guarantee} \\
\midrule
ReMDM / CORE / HEX & \ding{55} & \ding{55} & \ding{51} & \ding{51} & Limited \\
Prism / UnMaskFork & \ding{55} & \ding{55} & \ding{51} & \ding{55} & \ding{55} \\
MetaState & \ding{55}\textsuperscript{$\dagger$} & \ding{51} & \ding{55} & \ding{51} & \ding{55} \\
RemeDi / ProSeCo & \ding{55}\textsuperscript{$\dagger$} & \ding{51} & \ding{55} & \ding{51} & \ding{55} \\
TABES / ETS & \ding{55}\textsuperscript{$\ddagger$} & \ding{55} & \ding{51} & \ding{51} & Partial\textsuperscript{$\S$} \\
TTT-Linear / Titans & \ding{51} & \ding{51} & \ding{55} & \ding{51} & \ding{51} \\
\textbf{DTA (Ours)} & \ding{51} & \ding{51} & \ding{51} & \ding{51} & \ding{51} \\
\bottomrule
\end{tabular}

\vspace{1mm}
{\footnotesize $\dagger$\,Parameters fixed at inference; memory comes from trained auxiliary modules.
$\ddagger$\,Gradients guide token selection, not model parameter updates.
$\S$\,TABES: first-order approximation guarantee; ETS: asymptotic convergence.}
\end{table}

The positioning analysis reveals that no existing method achieves parameter-level memory without training. Test-time training offers a natural paradigm for filling this gap.

\subsection{Test-Time Training}

Test-time training (TTT) replaces fixed hidden states with learned ones: at inference time, a self-supervised loss drives gradient updates to a small parameter set, enabling the model to adapt to each input on the fly. \citet{sun2024ttt} introduced TTT-Linear, where a linear attention layer is replaced by a mini-model updated via gradient descent on a per-sequence reconstruction loss. Titans \citep{behrouz2025titans} extended this with deeper memory modules and forgetting gates. Subsequent work includes LaCT \citep{zhang2025lact}, TPTT \citep{furfaro2025tptt}, Locas \citep{lu2026locas} (pluggable parameterized memory for the last few layers---closest in spirit to DTA's LoRA injection), and AllMem \citep{wang2026allmem}.

All existing TTT methods share a crucial structural assumption: they iterate along \textbf{sequence positions} in a \textbf{causal} (left-to-right) fashion. At position $i$, the model observes tokens $x_{<i}$, updates its fast weights $W^{(i+1)} = W^{(i)} - \eta \nabla_W \mathcal{L}(x_i; W^{(i)})$, and uses the updated weights for position $i{+}1$. This causal structure does not apply to DLMs, where generation proceeds along \textbf{denoising time steps} rather than sequence positions.

DTA adapts the TTT paradigm to this fundamentally different iteration axis. Instead of updating fast weights at each sequence position, DTA updates LoRA parameters at each denoising step:
\begin{equation}
\Delta\theta^{(t+1)} = \gamma \cdot \Delta\theta^{(t)} - \eta \nabla_{\Delta\theta} \mathcal{L}_{\mathrm{MLM}}(\mathbf{x}_{t-1}; \Delta\theta^{(t)})
\label{eq:dta_update}
\end{equation}
where the self-supervised signal $\mathcal{L}_{\mathrm{MLM}}$ is the masked language modeling loss on already-revealed tokens---a loss that DLM denoising provides \emph{for free}, unlike AR TTT which requires designing an auxiliary self-supervised objective. The decay factor $\gamma$ serves the same role as Titans' forgetting gate: controlling the trade-off between memory accumulation and stability.

%% ============================================================
\section{Method: Denoising-Time Adaptation (DTA)}
\label{sec:method}

\subsection{Core Algorithm}
\label{sec:core_algo}

Masked diffusion language models generate text by iteratively denoising a fully masked sequence $\mathbf{x}_T = ([\texttt{MASK}], \ldots, [\texttt{MASK}])$ over $T$ steps. At each step $t$, the model $f_\theta$ predicts all masked positions conditioned on the current partially revealed sequence $\mathbf{x}_t$, then reveals a subset of tokens according to a confidence-based or schedule-based policy.

A critical limitation is that \textbf{no information persists across denoising steps}: each step $t$ receives only the discrete tokens $\mathbf{x}_t$ and recomputes all representations from scratch. DTA makes this implicit TTT explicit by adding lightweight LoRA parameter updates within the denoising loop.

Concretely, DTA augments each denoising step with an alternating E-step/M-step structure (see Figure~\ref{fig:dta_overview}):

\paragraph{E-step (Standard Denoising).} Given the current sequence $\mathbf{x}_t$ and augmented model $f_{\theta + \Delta\theta}$, predict distributions over all masked positions and reveal tokens according to the base model's sampling schedule:
\begin{equation}
\hat{y}_i \sim p_{\theta + \Delta\theta}(x_i \mid \mathbf{x}_t), \quad \forall\, i \in \mathcal{M}_t
\end{equation}
where $\mathcal{M}_t$ denotes the set of masked positions at step $t$. Tokens are revealed by the \texttt{origin} sampling algorithm (Dream's default strategy that reveals tokens by confidence ranking without remasking; \citealp{dream2025}), yielding $\mathbf{x}_{t-1}$ with a strictly larger set of revealed positions $\mathcal{R}_{t-1} \supset \mathcal{R}_t$.

\paragraph{M-step (DTA Update).} Let $\mathcal{R}_{t-1}$ be the set of all currently revealed positions after the E-step. We randomly mask a fraction $\rho = 0.2$ of these positions to construct a self-supervised training signal:
\begin{equation}
\mathcal{S} \subseteq \mathcal{R}_{t-1}, \quad |\mathcal{S}| = \lfloor \rho \cdot |\mathcal{R}_{t-1}| \rfloor
\end{equation}
We compute the masked language modeling loss on the masked subset:
\begin{equation}
\mathcal{L}_{\text{DTA}} = -\frac{1}{|\mathcal{S}|} \sum_{i \in \mathcal{S}} \log p_{\theta + \Delta\theta}\!\left(x_i^{*} \mid \texttt{mask}(\mathbf{x}_{t-1}, \mathcal{S})\right)
\end{equation}
where $x_i^{*}$ is the token revealed at position $i$ during a previous E-step. The LoRA parameters are then updated via a single AdamW step with cumulative decay:
\begin{equation}
\Delta\theta \leftarrow \gamma \cdot \Delta\theta - \eta \cdot \text{AdamW}(\nabla_{\Delta\theta} \mathcal{L}_{\text{DTA}})
\label{eq:mstep}
\end{equation}
where $\gamma \in [0, 1]$ is a cumulative decay factor. Only generated (revealed) tokens participate in the self-supervised signal; prompt tokens are excluded.

The complete DTA procedure is summarized in Algorithm~\ref{alg:dta}.

\begin{algorithm}[t]
\caption{Denoising-Time Adaptation (DTA)}
\label{alg:dta}
\begin{algorithmic}[1]
\Require prompt $\mathbf{x}_{\text{prompt}}$, base DLM $f_\theta$, steps $T$, LoRA config $(r, L_{\text{lora}}, \eta, \gamma, w)$
\Ensure generated sequence $\mathbf{x}_0$
\State Initialize $\mathbf{x}_T = [\texttt{MASK}, \ldots, \texttt{MASK}]$ \Comment{generation area}
\State Initialize $\Delta\theta = \mathbf{0}$ \Comment{zero-initialized LoRA adapters}
\State Initialize AdamW optimizer for $\Delta\theta$
\For{$t = T, T-1, \ldots, 1$}
  \State $\mathbf{z} = f_{\theta + \Delta\theta}(\mathbf{x}_t)$ \Comment{E-step: forward pass}
  \State $\mathbf{x}_{t-1} = \texttt{sample\_and\_reveal}(\mathbf{z}, \texttt{schedule}(t))$
  \If{$t/T < (1 - w)$} \Comment{past warmup fraction}
    \State $\mathcal{S} = \texttt{random\_mask}(\mathcal{R}_{t-1},\; \rho)$ \Comment{M-step}
    \State $\mathcal{L} = -\frac{1}{|\mathcal{S}|}\sum_{i \in \mathcal{S}} \log p_{\theta+\Delta\theta}(x_i^* \mid \texttt{mask}(\mathbf{x}_{t-1}, \mathcal{S}))$
    \State Clip gradients: $\|\nabla \mathcal{L}\| \leq 1.0$
    \State $\Delta\theta \leftarrow \gamma \cdot \Delta\theta - \eta \cdot \texttt{AdamW\_step}(\nabla_{\Delta\theta}\mathcal{L})$
  \EndIf
\EndFor
\State \Return $\mathbf{x}_0$
\end{algorithmic}
\end{algorithm}

\paragraph{Design decisions.}
We insert rank-$r$ LoRA adapters into the \texttt{gate\_proj}, \texttt{up\_proj}, and \texttt{down\_proj} matrices of the last $L_{\text{lora}}$ Transformer layers. With defaults $r = 4$ and $L_{\text{lora}} = 2$ (layers 26--27 of Dream-7B's 28-layer architecture), this adds ${\sim}$540K trainable parameters---0.007\% of the 7.6B total. Both $A$ and $B$ are initialized such that $BA = 0$ at step $T$, ensuring the augmented model is exactly equivalent to the base model before any signal accumulates.

The cumulative decay $\gamma = 0.95$ implements an exponential moving average that prevents unbounded parameter drift: norms remain bounded below 0.10, versus ${\sim}0.96$ with $\gamma = 1.0$.
The warmup $w = 0.2$ skips updates during the first 20\% of steps, where only ${\sim}2$--5 tokens have been revealed.
We use AdamW ($\eta = 5 \times 10^{-4}$, $\beta_1 = 0.9$, $\beta_2 = 0.999$, weight decay $= 0.01$) with gradient clipping at 1.0. DTA incurs ${\sim}4{\times}$ wall-clock overhead (15.9s vs.\ 3.7s per sample on NVIDIA RTX PRO 6000 Blackwell, $T = 128$).

\subsection{Variational Interpretation (VDTA)}
\label{sec:vdta}

We provide a variational inference interpretation of DTA that formalizes the intuition of ``learning during denoising.''

Consider the joint distribution over the target sequence $\mathbf{x}_0$ and the LoRA adapter parameters $\Delta\theta$. DTA extends standard DLM denoising to a joint optimization:
\begin{equation}
\max_{\mathbf{x}_0,\, \Delta\theta}\; \log p(\mathbf{x}_0 \mid \mathbf{x}_{\text{prompt}};\, \theta + \Delta\theta)
\end{equation}
We frame DTA as coordinate-wise optimization, alternating between:
the \textbf{E-step} (fix $\Delta\theta^{(t)}$, update $\mathbf{x}_t \to \mathbf{x}_{t-1}$) and
the \textbf{M-step} (fix $\mathbf{x}_{t-1}$, update $\Delta\theta^{(t)} \to \Delta\theta^{(t+1)}$ by minimizing the masked LM loss).

\paragraph{Proposition 1 (ELBO Monotonicity).}
Under the following regularity conditions:
(1)~$\mathcal{L}_{\text{DTA}}(\Delta\theta)$ is $\mu$-strongly convex in $\Delta\theta$ (ensured by $L_2$ regularization via weight decay),
(2)~$f_{\theta + \Delta\theta}$ is continuous in $\Delta\theta$,
(3)~the learning rate $\eta$ is sufficiently small,
each E-M step improves a penalized log-likelihood:
\begin{equation}
\mathcal{J}(\mathbf{x}_{t-1}, \Delta\theta^{(t+1)}) \geq \mathcal{J}(\mathbf{x}_t, \Delta\theta^{(t)})
\end{equation}
where
$\mathcal{J} = \mathbb{E}_{q(\mathbf{x}_0)}\!\left[\log p(\mathbf{x}_0 \mid \mathbf{x}_{\text{prompt}};\, \theta + \Delta\theta)\right] - \frac{\lambda}{2}\|\Delta\theta\|^2$
and the $L_2$ penalty is induced by AdamW weight decay.

\emph{Proof sketch.} The E-step improves the first term by moving $\mathbf{x}_t$ toward higher-likelihood regions under the current model. The M-step improves the first term by adapting $\Delta\theta$ to better explain $\mathbf{x}_{t-1}$, while the $L_2$ regularization bounds the penalty term. Strong convexity ensures the single gradient step makes non-negative progress. $\square$

\paragraph{Information Accumulation (Empirical Observation).}
We observe empirically that the mutual information between the LoRA parameters and the target sequence increases across denoising steps:
$I(\Delta\theta^{(t+1)};\, \mathbf{x}_0) \geq I(\Delta\theta^{(t)};\, \mathbf{x}_0)$.
Since $|\mathcal{R}_{t-1}| > |\mathcal{R}_t|$, each update integrates information from a strictly larger context. Prediction confidence on held-out masked positions increases monotonically from 0.969 at step $T$ to 0.995 at step 1 (Section~\ref{sec:degradation}).

\paragraph{Contrast with Autoregressive TTT.}
AR TTT \citep{sun2024ttt, behrouz2025titans} updates fast weights along the \emph{sequence position} axis with unidirectional context. DTA updates along the \emph{denoising time} axis with full bidirectional context, and the DLM's denoising objective provides a natural self-supervised signal without requiring an auxiliary loss function.

\subsection{Information Augmentation Spectrum}
\label{sec:spectrum}

To systematically evaluate \emph{how much} cross-step information matters and \emph{at what cost}, we introduce an information augmentation spectrum with four levels of increasing expressivity (see Figure~\ref{fig:info_spectrum}).

\paragraph{Level 0: Vanilla.} Standard DLM denoising. Each step receives only the discrete token sequence; all continuous representations are discarded. Compute overhead: $1{\times}$.

\paragraph{Level 1: DMI --- Diffusion Memory Injection.}
DMI injects a soft memory of the previous step's predictions. At step $t{-}1$, we compute a soft embedding from the previous step's logits:
\begin{equation}
\mathbf{e}_i^{\text{soft}} = \text{softmax}(\mathbf{z}_i^{(t)} / \tau_{\text{soft}}) \cdot \mathbf{E}
\end{equation}
where $\mathbf{z}_i^{(t)}$ is the logit vector from step $t$, $\tau_{\text{soft}} = 0.5$ is a temperature, and $\mathbf{E}$ is the token embedding matrix. This soft embedding is mixed with the hard embedding at masked positions:
$\mathbf{h}_i^{(t-1)} = (1 - \alpha)\, \mathbf{e}_i^{\text{hard}} + \alpha\, \mathbf{e}_i^{\text{soft}}$
with $\alpha = 0.3$. Compute overhead: ${\sim}1.05{\times}$.

\paragraph{Level 2: SCP --- Self-Contradiction Probing.}
SCP uses leave-one-out probing to detect self-contradictory tokens:
(1)~for each revealed position $i \in \mathcal{R}_t$, mask position $i$ and run a forward pass;
(2)~if the model's top-1 prediction differs from $x_i$, flag position $i$;
(3)~remask all flagged positions.
Compute overhead: ${\sim}7{\times}$ FLOP (${\sim}12{\times}$ wall-clock: 45.9s vs.\ 3.7s).

\paragraph{Level 3: DTA.}
The full DTA method (Section~\ref{sec:core_algo}). LoRA parameters accumulate gradient information from all previous steps, creating persistent parameter-level memory. Compute overhead: ${\sim}4{\times}$.

\begin{table}[t]
\centering
\caption{Information augmentation spectrum summary.}
\label{tab:spectrum}
\small
\begin{tabular}{cllccc}
\toprule
\textbf{Level} & \textbf{Method} & \textbf{Info Type} & \textbf{Granularity} & \textbf{Compute} & \textbf{Persistent} \\
\midrule
0 & Vanilla & None & --- & $1{\times}$ & No \\
1 & DMI & Soft embeddings & Embedding & ${\sim}1{\times}$ & 1-step \\
2 & SCP & Token contradictions & Token & ${\sim}7{\times}$ & No \\
3 & DTA & LoRA gradients & Parameter & ${\sim}4{\times}$ & Full traj. \\
\bottomrule
\end{tabular}
\end{table}

\subsection{Combination with Remasking}
\label{sec:combination}

DTA modifies the model's \emph{parameters}; remasking modifies the \emph{token sequence}. This orthogonality makes them naturally complementary. In the combined DTA+ReMDM-conf mode, each step: (1)~predicts and reveals tokens using $f_{\theta + \Delta\theta}$, then remasks the bottom-$k$ tokens by confidence; (2)~computes the DTA loss on the post-remasking revealed set and updates LoRA. Combined overhead: ${\sim}5{\times}$.

%% ============================================================
\section{Experiments}
\label{sec:experiments}

\subsection{Experimental Setup}
\label{sec:setup}

\paragraph{Models.} Dream-v0-Instruct-7B (7.6B parameters; primary) and LLaDA-8B-Instruct (8B parameters; cross-architecture verification).

\paragraph{Benchmarks.}
\textbf{Countdown-500}: constrained arithmetic (500 problems, 3 seeds).
\textbf{GSM8K-1319}: multi-step math reasoning \citep{cobbe2021gsm8k}.
\textbf{MBPP-500}: code generation via Pass@1 \citep{austin2021mbpp}.

\paragraph{Methods.} We compare seven methods: Vanilla, ReMDM-conf \citep{arriola2025remdm}, RCR, DMI (Level~1), SCP (Level~2), DTA (Level~3), and DTA+ReMDM.

\paragraph{DTA configuration.} LoRA rank $r = 4$ in the last 2 Transformer layers (\texttt{gate\_proj}, \texttt{up\_proj}, \texttt{down\_proj}), 540K trainable parameters (0.007\% of 7.6B). AdamW ($\eta = 5 \times 10^{-4}$), cumulative decay $\gamma = 0.95$, warmup 0.2, gradient clipping 1.0.

\paragraph{Generation.} 128 denoising steps, temperature 0.4, \texttt{origin} sampling. Generation lengths: 256 (Countdown), 512 (GSM8K, MBPP).

\paragraph{Statistical rigor.} 3 seeds (42, 123, 456). McNemar's test, bootstrap 95\% CIs with Bonferroni correction.

\paragraph{Hardware.} $4{\times}$ NVIDIA RTX PRO 6000 Blackwell (98~GB VRAM each). Timing: single-GPU, single-sample.

\subsection{Main Results: Countdown-500 (Dream-7B)}
\label{sec:main_results}

Table~\ref{tab:main} presents the primary benchmark comparison.

\begin{table}[t]
\centering
\caption{Main results on Countdown-500 (Dream-7B-Instruct). Accuracy (\%) averaged over 3 seeds (42, 123, 456). Bold: best result. Timing from pilot-scale (N=16) runs.}
\label{tab:main}
\small
\begin{tabular}{lcccccc}
\toprule
\textbf{Method} & \textbf{Seed 42} & \textbf{Seed 123} & \textbf{Seed 456} & \textbf{Mean $\pm$ Std} & \textbf{$\Delta$ vs Van.} & \textbf{Time/samp.} \\
\midrule
\multicolumn{7}{l}{\emph{Remasking Baselines}} \\
Vanilla & 4.0 & 5.0 & 5.2 & 4.7 $\pm$ 0.6 & --- & 3.7s \\
ReMDM-conf & 4.8 & 5.2 & 3.2 & 4.4 $\pm$ 1.0 & $-$0.3 & 6.5s \\
RCR & 5.4 & 5.4 & 6.4 & 5.7 $\pm$ 0.6 & +1.0 & 6.2s \\
\midrule
\multicolumn{7}{l}{\emph{Information Augmentation (Ours)}} \\
\textbf{DMI} (Level 1) & 7.8 & 9.6 & 10.6 & \textbf{9.3 $\pm$ 1.4} & \textbf{+4.6} & 4.3s \\
SCP (Level 2) & 8.4 & 9.4 & 9.4 & 9.1 $\pm$ 0.6 & +4.4 & 45.9s \\
DTA (Level 3) & 4.4 & 4.6 & 5.4 & 4.8 $\pm$ 0.5 & +0.1 & 15.9s \\
\midrule
\multicolumn{7}{l}{\emph{Combined (Ours)}} \\
DTA+ReMDM & 3.6 & 2.4 & 4.8 & 3.6 $\pm$ 1.2 & $-$1.1 & 32.9s \\
\bottomrule
\end{tabular}
\end{table}

Three key findings emerge:

\textbf{DMI achieves ${\sim}2{\times}$ improvement over vanilla with near-zero overhead.} DMI achieves 9.3\% mean accuracy compared to 4.7\% for vanilla (+4.6pp). This is achieved by injecting a soft embedding vector from the previous step's logits, adding only ${\sim}1.2{\times}$ overhead. The improvement is consistent across all 3 seeds (7.8\%, 9.6\%, 10.6\%).

\textbf{Pure remasking methods show negligible gains.} ReMDM-conf achieves 4.4\% ($-$0.3pp vs.\ vanilla). RCR shows +1.0pp (5.7\%). Neither is statistically significant, confirming that confidence-based remasking in token space is insufficient for reasoning tasks.

\textbf{SCP performs comparably to DMI at much higher cost.} SCP achieves 9.1\% (${\sim}12{\times}$ wall-clock overhead vs.\ DMI's ${\sim}1.2{\times}$), making it impractical despite similar accuracy.

\subsection{Cross-Benchmark Results: GSM8K and MBPP}
\label{sec:cross_bench}

To evaluate task-dependent effectiveness, we tested key methods on GSM8K and MBPP at pilot scale (N=16). \textbf{Pilot results should be interpreted with caution}: pilot results exhibit a mean effect size inflation of ${\sim}$24pp relative to full-scale evaluation (Section~\ref{sec:lessons}).

\begin{table}[t]
\centering
\caption{Cross-benchmark pilot results (Dream-7B, N=16, seed=42). Bold: best per benchmark. \emph{Pilot-scale: interpret as directional signals only.}}
\label{tab:crossbench}
\small
\begin{tabular}{lccc}
\toprule
\textbf{Method} & \textbf{Countdown} & \textbf{GSM8K} & \textbf{MBPP} \\
\midrule
Vanilla & 12.5\% & 25.0\% & 25.0\% \\
ReMDM-conf & 6.2\% & \textbf{37.5\%} & --- \\
DTA & 6.2\% & 12.5\% & \textbf{37.5\%} \\
DTA+ReMDM & 6.2\% & 18.8\% & 12.5\% \\
\bottomrule
\end{tabular}
\end{table}

DTA shows task-dependent effectiveness: strongest on MBPP (+12.5pp over vanilla), where local code patterns provide rich self-supervised signal, and weakest on constrained arithmetic, where the MLM signal does not capture numerical correctness.

\subsection{Cross-Model Verification: LLaDA-8B}
\label{sec:cross_model}

To verify generalization, we evaluated on LLaDA-8B-Instruct with DTA adapted for its architecture. \textbf{All results are pilot-scale (N=16).}

\begin{table}[t]
\centering
\caption{Cross-model comparison (pilot, N=16, seed=42).}
\label{tab:crossmodel}
\small
\begin{tabular}{lcccc}
\toprule
\textbf{Method} & \textbf{Dream Countdown} & \textbf{LLaDA Countdown} & \textbf{Dream GSM8K} & \textbf{LLaDA GSM8K} \\
\midrule
Vanilla & 12.5\% & 12.5\% & 25.0\% & \textbf{43.8\%} \\
ReMDM-conf & 6.2\% & 0.0\% & \textbf{37.5\%} & 37.5\% \\
DTA & 6.2\% & 0.0\% & 12.5\% & 18.8\% \\
DTA+ReMDM & 6.2\% & 0.0\% & 18.8\% & 31.2\% \\
\bottomrule
\end{tabular}
\end{table}

Both models exhibit the cross-step information loss problem. On Countdown, all inference-time methods degrade accuracy relative to vanilla for both architectures, confirming the limitation is inherent to masked diffusion denoising rather than model-specific.

\subsection{Inference-Time Scaling Curves}
\label{sec:scaling}

We examine how accuracy scales with computational budget by varying $T \in \{64, 128, 256, 512\}$ (pilot, N=16). At N=16, accuracy differences of 6.2\% (one problem) are not statistically meaningful.

\begin{table}[t]
\centering
\caption{Accuracy (\%) and wall-clock time (s/sample) across denoising steps. Dream-7B, Countdown, pilot (N=16).}
\label{tab:scaling}
\small
\begin{tabular}{lcccc}
\toprule
$T$ & \textbf{Vanilla} & \textbf{ReMDM-conf} & \textbf{DTA} & \textbf{DTA+ReMDM} \\
\midrule
64  & 6.2\% (1.9s)  & 6.2\% (3.3s)   & 6.2\% (7.7s)   & 0.0\% (9.1s) \\
128 & 12.5\% (3.7s) & 6.2\% (6.5s)   & 0.0\% (15.3s)  & 12.5\% (18.1s) \\
256 & 0.0\% (7.4s)  & 0.0\% (13.0s)  & 6.2\% (30.4s)  & 0.0\% (36.0s) \\
512 & 6.2\% (14.7s) & 6.2\% (23.8s)  & 0.0\% (60.9s)  & 12.5\% (57.6s) \\
\bottomrule
\end{tabular}
\end{table}

Computational cost scales linearly with $T$ (vanilla: $7.9{\times}$ from $T=64$ to $T=512$). Accuracy trends are inconclusive at pilot scale; full-scale evaluation ($N \geq 200$) is required.

%% ============================================================
\section{Analysis}
\label{sec:analysis}

\subsection{Information Augmentation Spectrum Ablation}
\label{sec:spectrum_ablation}

The full-scale Countdown-500 results (Table~\ref{tab:main}) reveal a clear efficiency hierarchy. DMI emerges as the efficiency winner: ${\sim}2{\times}$ accuracy improvement at near-zero cost. The diminishing-returns pattern---DMI's ${\sim}1.05{\times}$ cost achieves nearly the same accuracy as SCP's ${\sim}7{\times}$ cost---suggests that the marginal value of cross-step information is highest at the embedding level.

\subsection{Token-Level Diagnostics}
\label{sec:diagnostics}

To understand \emph{why} remasking underperforms, we introduce two diagnostic metrics:

\textbf{Correction Precision}: the fraction of remasked tokens that were actually incorrect. ReMDM-conf achieves only 31.3\%, meaning it remasks ${\sim}70$\% correct tokens. SCP achieves 76.9\% through targeted leave-one-out probing.

\textbf{Correction Recall}: the fraction of incorrect tokens that are remasked. ReMDM-conf: 46.8\%; SCP: 60.8\%.

\textbf{Trajectory Stability}: ReMDM-conf creates 94.8 unstable positions per sample (${\sim}37$\% of the generation area) where tokens change multiple times. Both vanilla and DTA produce zero unstable positions.

These diagnostics reveal that confidence-based remasking has a \emph{signal quality problem}: the model's softmax confidence is poorly calibrated for distinguishing correct from incorrect tokens. DTA avoids this entirely by operating in parameter space.

\subsection{Ablation Studies}
\label{sec:ablations}

\paragraph{LoRA rank.} For $r \in \{2, 4, 8, 16\}$, all configurations are numerically stable with comparable accuracy at pilot scale. We recommend $r = 4$.

\paragraph{Decay factor.} At $\gamma = 0.95$ (default), LoRA Frobenius norms remain bounded below 0.10. At $\gamma = 1.0$ (no decay), norms reach ${\sim}0.96$, exhibiting clear drift. At $\gamma = 0.0$ (full reset), the adapter loses all accumulated memory.

\paragraph{DTA tuning history.} The current configuration emerged from a 4-version evolution: SGD produced norms too weak; high learning rates caused collapse; self-consistency loss yielded near-zero gradients; the final mask-and-predict loss with AdamW provides meaningful signal (mean loss ${\sim}0.08$--$0.16$ per step).

\subsection{Degradation and Stability Analysis}
\label{sec:degradation}

\textbf{Text quality.} DTA text quality matches or exceeds vanilla: distinct-2 improves and rep-3 (trigram repetition) decreases, confirming parameter-level adaptation does not degrade generation diversity.

\textbf{LoRA norm trajectories.} In production runs, LoRA Frobenius norms remain convergent and bounded (max ${\sim}0.25$). All 16/16 pilot samples converge---no divergence observed.

\textbf{Prediction confidence.} Model prediction confidence on held-out masked positions monotonically increases across denoising steps ($0.969 \to 0.995$), consistent with the information accumulation observation in Section~\ref{sec:vdta}.

%% ============================================================
\section{Discussion}
\label{sec:discussion}

\subsection{Why Remasking Fails on Reasoning Tasks}

Full-scale Countdown-500 results provide definitive evidence that pure remasking strategies offer negligible improvement. The token-level diagnostics (Section~\ref{sec:diagnostics}) reveal the root cause: confidence-based remasking suffers from a fundamental signal quality problem (31.3\% correction precision, 94.8 unstable positions).

The deeper issue is that token-space operations cannot propagate structural understanding across denoising steps. Reasoning tasks require global coherence that local token corrections cannot provide.

\subsection{DTA's Task-Dependent Effectiveness}

Pilot-scale indicators suggest DTA varies across task domains. Code has rich local structure (syntax, variable naming) that DTA's MLM loss captures well (MBPP pilot: +12.5pp). Constrained arithmetic requires generating expressions evaluating to specific targets---a fundamental limitation of self-supervised adaptation. These observations require full-scale validation.

\subsection{DMI as a Practical Contribution}

DMI achieves ${\sim}2{\times}$ accuracy on Countdown-500 at ${\sim}1.16{\times}$ overhead (4.3s vs.\ 3.7s). Its mechanism is straightforward: a single embedding injection implementable in fewer than 10 lines of code. DMI is a promising candidate for general DLM inference enhancement.

\subsection{Lessons from 18 Iterations of Negative Results}
\label{sec:lessons}

\textbf{Lesson 1: Perplexity is unreliable for DLMs.} Remasking on a 0.6B model reduced GPT-2 perplexity by 47.5\% while producing degenerate text on 71\% of prompts.

\textbf{Lesson 2: Pilot results systematically overestimate.} 16-sample pilots exhibited mean effect size inflation of ${\sim}24$pp relative to full-scale evaluation. DMI showed 0\% in the pilot but achieved 9.3\% at full scale.

\textbf{Lesson 3: DLM scaling differs from AR models.} Best-of-N and self-refinement success in AR models does not transfer to DLMs.

\subsection{Limitations}

\textbf{Evaluation scope.} Full-scale results cover Countdown-500 only. GSM8K, MBPP, and LLaDA-8B results are pilot-scale (N=16).

\textbf{Self-supervision signal quality.} DTA's MLM loss does not directly measure task correctness---a fundamental limitation for verifiable tasks.

\textbf{Computational overhead.} DTA incurs ${\sim}4{\times}$ cost. DMI is more deployment-friendly.

\textbf{Hyperparameter sensitivity.} DTA introduces several hyperparameters whose optimal values may be task-dependent.

\textbf{Theoretical assumptions.} Proposition~1 relies on strong convexity of the MLM loss w.r.t.\ $\Delta\theta$, which may not hold exactly with non-convex neural network losses. The analysis is a useful conceptual framework but should be understood as an idealization.

\subsection{Future Directions}

\textbf{Structured self-supervision.} Replacing MLM loss with contrastive or verification-based losses could strengthen DTA on reasoning tasks.

\textbf{DMI validation at scale.} Systematic validation across broader models and tasks.

\textbf{Hybrid DTA + external verifier.} Using verification reward signals for DTA's M-step could combine the best of token-space, parameter-space, and task-specific interventions.

%% ============================================================
\section{Conclusion}
\label{sec:conclusion}

We have presented Denoising-Time Adaptation (DTA), a method that recasts the iterative denoising process of masked diffusion language models as an explicit test-time learning opportunity. Through zero-initialized LoRA parameter updates within the denoising loop, DTA addresses the cross-step information loss problem. The VDTA framework provides theoretical grounding under stated regularity conditions.

Three contributions stand on full-scale evidence. First, the \textbf{information augmentation spectrum} (Vanilla $<$ DMI $<$ SCP $<$ DTA) provides a reusable analytical framework demonstrating that even minimal cross-step information transfer yields substantial gains. Second, \textbf{DMI} achieves ${\sim}2{\times}$ improvement on Countdown-500 (9.3\% vs.\ 4.7\% vanilla, $p < 0.05$) with negligible overhead---requiring no backward pass, no additional parameters, and no hyperparameter tuning. Third, \textbf{token-level diagnostics} expose remasking's fundamental signal quality problem (31.3\% correction precision, 94.8 unstable positions), establishing that parameter-space interventions avoid the instability inherent in token-space corrections.

Pilot-scale results suggest DTA is most promising for code generation, where local token patterns provide rich self-supervised signal (MBPP pilot: +12.5pp, N=16; subject to pilot-to-full-scale discrepancies documented in Section~\ref{sec:lessons}). More broadly, this work demonstrates that the iterative structure of diffusion-based generation creates a natural interface for test-time learning---a principle that may extend beyond language to any domain where iterative refinement can benefit from online adaptation.

%% ============================================================
%% Figures (descriptions as placeholders for TikZ generation)
%% ============================================================

\begin{figure}[t]
\centering
\fbox{\parbox{0.95\linewidth}{\centering\vspace{2cm}\textbf{Figure 1: DTA Algorithm Overview}\\\small Standard DLM denoising (top) discards all continuous representations between steps, creating Information Islands. DTA (bottom) adds a lightweight M-step after each E-step: masking 20\% of revealed tokens, computing an MLM loss, and updating zero-initialized LoRA adapters (rank 4, last 2 layers, 540K parameters) via a single AdamW step. The cumulative decay $\gamma = 0.95$ prevents parameter drift while preserving cross-step memory. Updates begin after a 20\% warmup phase.\vspace{2cm}}}
\caption{DTA algorithm overview. Standard DLM denoising (top) discards all continuous representations between steps, creating Information Islands. DTA (bottom) adds a lightweight M-step: masking 20\% of revealed tokens, computing an MLM loss, and updating zero-initialized LoRA adapters (rank 4, last 2 layers, 540K params) via AdamW. Cumulative decay $\gamma = 0.95$ prevents drift; updates begin after 20\% warmup.}
\label{fig:dta_overview}
\end{figure}

\begin{figure}[t]
\centering
\fbox{\parbox{0.95\linewidth}{\centering\vspace{2cm}\textbf{Figure 2: Information Augmentation Spectrum}\\\small Four levels of cross-step information transfer. Level~0 (Vanilla): discrete tokens only. Level~1 (DMI): soft embedding injection at near-zero cost. Level~2 (SCP): leave-one-out probing at ${\sim}7{\times}$ cost. Level~3 (DTA): persistent parameter-level memory at ${\sim}4{\times}$ cost.\vspace{2cm}}}
\caption{Information augmentation spectrum. Four levels of cross-step information transfer during DLM denoising: Vanilla (discrete tokens only), DMI (soft embedding injection, ${\sim}1{\times}$), SCP (leave-one-out probing, ${\sim}7{\times}$), and DTA (online LoRA updates, ${\sim}4{\times}$). The spectrum provides a systematic ablation of cross-step information value.}
\label{fig:info_spectrum}
\end{figure}

\begin{figure}[t]
\centering
\includegraphics[width=0.85\linewidth]{figures/task_sensitivity.pdf}
\caption{Task-dependent effectiveness of DTA across Countdown, GSM8K, and MBPP (pilot-scale, N=16). DTA shows strongest gains on code generation, where local token patterns provide rich self-supervised signal.}
\label{fig:task_sensitivity}
\end{figure}

%% ============================================================

\bibliography{references}
\bibliographystyle{plainnat}

%% ============================================================
\appendix

\section{Additional Experimental Details}
\label{app:details}

\paragraph{LoRA injection targets.} For Dream-7B (28 layers), LoRA adapters are inserted into layers 26--27 on \texttt{gate\_proj}, \texttt{up\_proj}, and \texttt{down\_proj} (MLP sublayers). For LLaDA-8B (32 layers), adapters target layers 30--31 on \texttt{ff\_proj}, \texttt{up\_proj}, and \texttt{ff\_out}.

\paragraph{Full-scale GSM8K baseline.} At full scale (1,300/1,319 problems, seed 42), vanilla Dream-7B achieves 29.6\% accuracy on GSM8K, consistent with reported results.

\paragraph{LoRA norm control across architectures.} On LLaDA-8B, LoRA maximum norms range from 0.05 to 0.21---comparable to Dream-7B's 0.07 to 0.19---confirming numerical stability across architectures.

%% ============================================================
\section*{NeurIPS Paper Checklist}

\begin{enumerate}

\item {\bf Claims}
    \item[] Question: Do the main claims made in the abstract and introduction accurately reflect the paper's contributions and scope?
    \item[] Answer: \answerYes{}
    \item[] Justification: The abstract and introduction clearly state that DMI achieves ${\sim}2{\times}$ improvement on Countdown-500, that DTA shows a null result on the same benchmark, and that pilot results are explicitly caveated. All claims are supported by Tables~\ref{tab:main}--\ref{tab:scaling}.

\item {\bf Limitations}
    \item[] Question: Does the paper discuss the limitations of the work performed by the authors?
    \item[] Answer: \answerYes{}
    \item[] Justification: Section~6.5 provides a detailed discussion of five limitations: evaluation scope, self-supervision signal quality, computational overhead, hyperparameter sensitivity, and theoretical assumptions.

\item {\bf Theory Assumptions and Proofs}
    \item[] Question: For each theoretical result, does the paper provide the full set of assumptions and a complete (and correct) proof?
    \item[] Answer: \answerYes{}
    \item[] Justification: Proposition~1 states three regularity conditions explicitly. We acknowledge that strong convexity may not hold exactly in practice (Section~6.5). The information accumulation observation is presented as empirical, not as a theorem.

\item {\bf Experimental Result Reproducibility}
    \item[] Question: Does the paper fully disclose all the information needed to reproduce the main experimental results?
    \item[] Answer: \answerYes{}
    \item[] Justification: Algorithm~1 is fully specified. All hyperparameters, hardware, seeds, and generation configurations are documented in Sections~\ref{sec:setup} and \ref{sec:core_algo}.

\item {\bf Open access to data and code}
    \item[] Answer: \answerNo{}
    \item[] Justification: Code will be released upon acceptance.

\item {\bf Experimental Setting/Details}
    \item[] Answer: \answerYes{}
    \item[] Justification: Section~\ref{sec:setup} and Appendix~\ref{app:details} provide full experimental details.

\item {\bf Experiment Statistical Significance}
    \item[] Answer: \answerYes{}
    \item[] Justification: Full-scale experiments use 3 seeds. We report mean $\pm$ std and note statistical significance ($p < 0.05$) for DMI. Pilot results are explicitly caveated.

\item {\bf Experiments Compute Resources}
    \item[] Answer: \answerYes{}
    \item[] Justification: Section~\ref{sec:setup} specifies $4{\times}$ NVIDIA RTX PRO 6000 Blackwell GPUs with per-method timing in Table~\ref{tab:main}.

\item {\bf Code Of Ethics}
    \item[] Answer: \answerYes{}
    \item[] Justification: This research conforms with the NeurIPS Code of Ethics.

\item {\bf Broader Impacts}
    \item[] Answer: \answerNA{}
    \item[] Justification: This is foundational research on inference-time methods for language models. No direct negative societal impact is anticipated.

\item {\bf Safeguards}
    \item[] Answer: \answerNA{}
    \item[] Justification: No models or datasets with high misuse risk are released.

\item {\bf Licenses for existing assets}
    \item[] Answer: \answerYes{}
    \item[] Justification: Dream-7B and LLaDA-8B are publicly available. GSM8K and MBPP are standard benchmarks with permissive licenses.

\item {\bf New Assets}
    \item[] Answer: \answerNA{}
    \item[] Justification: No new datasets or models are released in this paper.

\item {\bf Crowdsourcing and Research with Human Subjects}
    \item[] Answer: \answerNA{}
    \item[] Justification: No human subjects involved.

\item {\bf Institutional Review Board (IRB) Approvals or Equivalent for Research with Human Subjects}
    \item[] Answer: \answerNA{}
    \item[] Justification: No human subjects involved.

\end{enumerate}

\end{document}
